Coupler graph of IQM Garnet and the couplings that the LRB-D circuits need. Garnet has 20 qubits on a square lattice with 30 tunable couplers~\cite{IQM240812433}, which join the only pairs on which its two-qubit gate $CZ$ acts (\texttt{qubit\_connectivity} of the \texttt{fake\_garnet} backend in IQM's \texttt{qiskit-on-iqm}). In (a)--(c) each qubit sits at its lattice site, so couplers are nearest-neighbor bonds, and empty sites show QB indices. Every coupler joins the two sublattices (white and gray), so every closed loop of couplers has even length. A placement puts a circuit's qutrits on distinct qubits so that a coupler joins every pair acted on by a two-qutrit gate (placements related by a symmetry count separately). A circuit with a placement fits Garnet. Purple marks $q_2$--$q_3$, used by the logical gates, and $a_j$ is the syndrome ancilla of $G_j$. Marks below each panel state whether each requirement is met.
(a)~Terminal postselection (const-0) needs the couplings $q_2$--$q_1$, $q_1$--$q_0$, $q_3$--$q_0$ and $q_4$--$q_2$ of the encoder of $\ket{\bar +}$ and $q_2$--$q_3$ for $CZ_{2\to 3}$ in $\bar S$ and $\mathrm{SWAP}_{2,3}$ in $\bar H$. It fits in 128 placements (one shown).
(b)~The one-ancilla check round alone fits in 32 placements, each the distance-2 surface-code patch with $q_1$, the qutrit in all four stabilizers, at the center. In all of them $q_2$ and $q_3$ are four couplers apart, so $\bar S$ and $\bar H$ would need SWAP routing. Panels (b)--(e) omit the encoder of $\ket{\bar +}$, which the simulation applies without noise and whose data--data couplings no placement in (b) or (c) provides.
(c)~The intended two-ancilla layout: the patch as the paper draws it, a conventional surface-code grid with ancilla sites N, W, E and S between the data. No site touches all three data qutrits of $G_1$ or of $G_4$, so each needs two sites. The panel uses N ($q_0$, $q_1$) with W ($q_3$) for $G_1$ and E ($q_1$, $q_4$) with S ($q_3$) for $G_4$, so N and E, which also measure $G_3$ and $G_2$, serve two stabilizers each. The sites share one sublattice, so no coupler joins N to W or E to S (dashed), where a relay would need its direct ancilla--ancilla gate. Here too $q_2$ and $q_3$ are two couplers apart. The simulated relay circuits use such direct gates, six dedicated ancillas instead of four sites, and a $G_4$ split that puts $q_1$ and $q_3$ on one ancilla, a pair that no site touches.
(d)~Couplings of the one-ancilla check round plus $q_2$--$q_3$, drawn as a graph with data qutrits gray and ancillas white as in (b). Gray rims mark qutrits off the loop. Every check coupling joins a gray and a white qutrit, but the crossed-out coupling $q_2$--$q_3$ joins two gray ones and closes the five-coupler loop $a_1$--$q_0$--$a_3$--$q_2$--$q_3$. Two shades cannot alternate around a loop of odd length, so the circuit fits no square lattice.
(e)~Couplings of the simulated two-ancilla check round, shaded in the same way. The crossed-out couplings $a_1$--$r_1$ and $a_4$--$r_4$ join two white qutrits, a syndrome ancilla and its relay ancilla $r_j$. Each closes a five-coupler loop, $a_1$--$q_1$--$a_4$--$q_3$--$r_1$ or $a_4$--$q_1$--$a_2$--$q_4$--$r_4$. Either loop alone rules out a square lattice, even without the logical gates.
